% !TEX TS-program = lualatex+se
% !TEX encoding = UTF-8

\documentclass[12pt, a4paper]{extbook}
% --- USTAWIENIA STRONY ---
\usepackage{geometry}
\geometry{
	%showframe,
	left=15mm,
	right=12mm,
	top=12mm,
	bottom=15mm,
	marginparwidth=0mm,
	marginparsep=0mm
	}
% usuwanie nagłówków/stopek z\,dokumentu
\usepackage[indent=0pt]{parskip} % ustawianie paragrafów
%\setstretch{1.5} % półtorej linii odległości między wierszami
\usepackage{titlesec} % ustawianie nagłówków
\usepackage{fullpage} % redukcja marginesów
\usepackage{setspace} % odstępy między liniami
\usepackage{enumitem}% listy numerowane
% -----------------------------

% --- CZCIONKI DOKUMENTU ---
\usepackage[latin,polish]{babel}
\usepackage{fontspec}
\usepackage{times}
%\usepackage[math]{anttor}
%\usepackage{tgtermes}
%\usepackage{mathptmx}
\usepackage{ebgaramond}
\usepackage{libertine}
\defaultfontfeatures{Ligatures=TeX}

%\setmainfont{TeX Gyre Termes}
%\setmainfont{TeX Gyre Heros}
%\setmainfont{Calluna Regular}
\setmainfont{EB Garamond}
%\setmainfont{Bodoni Moda}
%\newfontfamily{\sekfont}{Times New Roman}
\newfontfamily{\sekfont}{Libertinus Serif}
\newcommand\sekcjafont{\fontsize{16}{16}\selectfont \sekfont}

% style nagłówków
\newcommand{\headingone}{\scshape\Large}
\newcommand{\headingtwo}{\large}
\newcommand{\margintext}[1]{\textit{\small #1}}
% przesunięcia poziome
\newcommand{\smallwidth}{\hskip 1em}
\newcommand{\medwidth}{\hskip 3pt}
\newcommand{\largewidth}{\hskip 12pt}
% przesunięcia pionowe
\newcommand{\linespace}{\vskip 4pt}
\newcommand{\paragraphspace}{\vskip 6pt}
\newcommand{\segmentspace}{\vskip 12pt}
\newcommand{\sectionspace}{\vskip 24pt}

%----- Strona tytułowa -------------
\newfontfamily{\tytfont}{JSL Ancient}
\newcommand\tytulfont[1]{\fontsize{36}{30}\selectfont \tytfont#1\par}
\newcommand\podtytulfont[1]{\fontsize{18}{10}\selectfont \tytfont#1\par}

% -----------------------------
\titlespacing*{\chapter}{0pt}{0pt}{5pt}	
\titleformat{\chapter}[hang]{\filcenter\LARGE\scshape\bfseries}{}{0pt}{} % format nagłówków
\titlespacing*{\section}{0pt}{0pt}{0pt}	
\titleformat{\section}[hang]{\filcenter\large\sekcjafont\scshape}{}{0pt}{} % format nagłówków

\newcommand{\sekcja}[2]{
	\begin{spacing}{0.6}
	\begin{center}
		\section*{#1}
		\begin{minipage}[c]{0.88\linewidth}
			\red{\small\textit{#2}}
		\end{minipage}
	\end{center}	
	\end{spacing}
	\vspace{4mm}}

% --- GRAFIKA I KOLORY ---
\usepackage{xcolor} % kolory
\usepackage{graphicx} % umożliwia import obrazów
\graphicspath{{img/}} % katalog obrazów
\usepackage{pgfornament}% ornamenty graficzne
% --- KONIEC GRAFIKA I KOLORY ----------------------

% --- KOMENDY WŁASNE ---
\newcommand{\rubryka}[1]{
	\begin{spacing}{0.8}
	\centering 
	\parbox{0.9\textwidth}{\centering\textcolor{red}{\small\textit{#1}}}
	\end{spacing}
	\vspace{3mm}
	\normalsize\normalfont}
%---------------------------------
\newcommand{\dekor}[2]{
	\begin{center}
		\vspace*{\fill}
		\includegraphics[width=#1\textwidth]{#2}
		\vspace*{\fill}
	\end{center}}	

% --- GREGORIO ---	
\usepackage{gregoriotex}
\gresetlinecolor{gregoriocolor}
\gresetheadercapture{annotation}{greannotation}{string}% odczytuje adnotacje nad inicjałem jako ciąg, co pozwala na jego renderowanie
\gresetheadercapture{commentary}{grecommentary}{string}% odczytuje komentarz jako ciąg, co pozwala na jego renderowanie
\grechangedim{commentaryraise}{0.3cm}{scalable} % podnosi komentarz wystarczająco ponad górną linię
\grechangestyle{annotation}{\color{black}\small}% format adnotacji
\grechangestyle{commentary}{\color{gregoriocolor}\small}% format komentarza
\grechangestyle{abovelinestext}{\color{gregoriocolor}\small}% format napisu nad liniami
\grechangestyle{initial}{\fontsize{36}{36}\selectfont} % zmniejsza wielkość inicjału do bardziej proporcjonalnego rozmiaru (=36pt)
\grechangedim{intersyllablespacenotes}{0.18cm}{scalable}% Minimalna odległość między nutami różnych sylab.

% umieszczenie zawartości składu ich unikalnego argumentu po lewej i prawej stronie każdej linii
\let\grelocalleftbox\localleftbox
\let\grelocalrightbox\localrightbox
\newcommand{\separator}{\vskip 6pt \centering\greseparator{3}{20}}

\usepackage{paracol}
\usepackage{xcolor} 
\setcolumnwidth{.49\textwidth, .49\textwidth}
\setlength{\columnseprule}{0.4pt}
\setlength{\columnsep}{.02\textwidth}
\colseprulecolor{red}
% ----------------------------
%------- Zarządzanie wersami w psalmach z tlumaczeniem -------------------------
\newcounter{licznik}
\newcommand{\licznikZero}{\setcounter{licznik}{0}}% zeruj licznik
\newcommand{\licznikUp}{\refstepcounter{licznik}}% zwiększ licznik 
\licznikZero
\newcommand{\BeginParallel}{\begin{spacing}{0.8} \begin{paracol}{2} \licznikZero}
% --- KOMENDY WŁASNE ---
\newcommand{\red}[1]{\textcolor{red}{#1}}
\newcommand{\black}[1]{\textcolor{black}{#1}}
% --------------------------------
\newcommand{\lat}[1]{\selectlanguage{latin}\ensurevspace{2\baselineskip}\begin{minipage}[t]{\linewidth}\ifnum \value{licznik}>0 \licznikUp \thelicznik .~#1 \else \licznikUp#1 \fi \end{minipage}\switchcolumn}
\newcommand{\latN}[1]{\selectlanguage{latin}\ensurevspace{2\baselineskip}\begin{minipage}[t]{\linewidth}#1\end{minipage}\switchcolumn}% kolumna bez numeracji wersów
\newcommand{\pol}[1]{\selectlanguage{polish}\begin{minipage}[t]{\linewidth}#1\end{minipage}\switchcolumn*}
\newcommand{\EndParallel}{\end{paracol} \end{spacing}}
\newcommand{\tlumacz}[1]{\textit{#1}}
\newcommand{\inicjal}[2]{\textcolor{red}{\LARGE\textbf{#1}}{#2}}



\begin{document}
\pagenumbering{gobble} % nie jest wyprowadzany żaden numer strony, chociaż numer jest nadal resetowany. Odniesienia do tej strony również są puste.
%\begin{center}
%	\includegraphics[width=1.0\textwidth]{Purification_BVM02.jpg}
%	\vfill
%	\podtytulfont{AD I. VESPERAS}
%	\tytulfont{In Purificatione\\Beatæ Mariæ Virginis}
%	\vfill
%	\large\textit{Liber usualis, Solesmes, 1961, p.1355}
%\end{center}
%\newpage
\pagenumbering{arabic}
\large
%% ----------------- Modlitwy wstępne -----------
\red{\large{Psalm 112}\linespace}
\gresetinitiallines{1}\greannotation{2. Ant.}\greannotation{3. a2}\gabcsnippet{(c4) <c>Q</c>Uan(e)do(e[ll:1]d) na(g)tus(hj) es(j'_) <c><v>\greheightstar</v></c> (,) in(j)ef(i')fa(g)bí(i')li(j)ter(h) ex(gf) Vír(g_[uh:l]h)gi(h)ne,(g.) (;) tunc(g) im(g)plé(i)tae(j) sunt(hgh___) Scri(gh)ptú(e.)rae :(e.) (:) sic(g)ut(h) plú(j')vi(i)a(jkj) in(i') vel(j)lus(h') de(h)scen(g')dí(i)sti,(h.) (;) ut(h) sal(j')vum(i) fá(j')ce(i)res(g) ge(hv_GF)nus(g) hu(gf)má(e.)num :(e.) (:) te(h_i) lau(g)dá(gf)mus(d_e) De(g)us(ghg) no(e.)ster.(e.) (::)}
\paragraphspace
\greannotation{Psalm}\greannotation{112}\gabcsnippet{(c4)<c>L</c>au(g)dá(hj)te,(j) <b>pú</b>(k)e(j)ri,(j) <b>Dó</b>(j)<b>mi</b>(ih)num:(j.) <c><v>\greheightstar</v></c>(:) lau(j)dá(j)te(j) <i>no</i>(ji)<i>men</i>(hi) <b>Dó</b>(h)mi(g)ni.(gh..) (::)
<c>2. </c><alt>fit reverentia</alt>Sit(j) no(j)men(j) Dó(j)mi(j)ni(j) <b>be</b>(k)ne(j)<b>díc</b>(ih)tum,(j.) <c><v>\greheightstar</v></c>(:) ex(j) hoc(j) nunc,(j) et(j) us(j)<i>que</i>(ji) <i>in</i>(hi) <b>sǽ</b>(h)cu(g)lum.(gh..) (::)
<c>3. </c>A(j) so(j)lis(j) or(j)tu(j) us(j)que(j) <b>ad</b>(k) oc(j)<b>cá</b>(ih)sum,(j.) <c><v>\greheightstar</v></c>(:) lau(j)dá(j)bi(j)le(j) <i>no</i>(ji)<i>men</i>(hi) <b>Dó</b>(h)mi(g)ni.(gh..) (::)
<c>4. </c>Ex(j)cél(j)sus(j) su(j)per(j) om(j)nes(j) <b>gen</b>(k)tes(j) <b>Dó</b>(j)<b>mi</b>(ih)nus,(j.) <c><v>\greheightstar</v></c>(:) et(j) su(j)per(j) cæ(j)los(j) gló(j)<i>ri</i>(ji)<i>a</i>(hi) <b>e</b>(h)jus.(gh..) (::)
<c>5. </c>Quis(j) sic(j)ut(j) Dó(j)mi(j)nus,(j) De(j)us(j) nos(j)ter,(j) qui(j) in(j) <b>al</b>(k)tis(j) <b>há</b>(j)<b>bi</b>(ih)tat,(j.) <c><v>\greheightstar</v></c>(:) et(j) hu(j)mí(j)li(j)a(j) ré(j)spi(j)cit(j) in(j) cæ(j)lo(j) <i>et</i>(ji) <i>in</i>(hi) <b>ter</b>(h)ra?(gh..) (::)
<c>6. </c>Sú(j)sci(j)tans(j) a(j) <b>ter</b>(k)ra(j) <b>ín</b>(j)<b>o</b>(ih)pem,(j.) <c><v>\greheightstar</v></c>(:) et(j) de(j) stér(j)co(j)re(j) é(j)<i>ri</i>(ji)<i>gens</i>(hi) <b>páu</b>(h)pe(g)rem:(gh..) (::)
<c>7. </c>Ut(j) cól(j)lo(j)cet(j) e(j)um(j) <b>cum</b>(k) prin(j)<b>cí</b>(j)<b>pi</b>(ih)bus,(j.) <c><v>\greheightstar</v></c>(:) cum(j) prin(j)cí(j)pi(j)bus(j) pó(j)<i>pu</i>(ji)<i>li</i>(hi) <b>su</b>(h)i.(gh..) (::)
<c>8. </c>Qui(j) ha(j)bi(j)tá(j)re(j) fa(j)cit(j) sté(j)ri(j)<b>lem</b>(k) in(j) <b>do</b>(ih)mo,(j.) <c><v>\greheightstar</v></c>(:) ma(j)trem(j) fi(j)li(j)ó(j)<i>rum</i>(ji) <i>læ</i>(hi)<b>tán</b>(h)tem.(gh..) (::)
<c>9. </c>Gló(j)ri(j)a(j) <b>Pa</b>(k)tri,(j) et(j) <b>Fí</b>(j)<b>li</b>(ih)o,(j.) <c><v>\greheightstar</v></c>(:) et(j) Spi(j)rí(j)<i>tu</i>(ji)<i>i</i>(hi) <b>Sanc</b>(h)to.(gh..) (::)
<c>10. </c>Sic(j)ut(j) e(j)rat(j) in(j) prin(j)cí(j)pi(j)o,(j) et(j) <b>nunc</b>,(k) et(j) <b>sem</b>(ih)per,(j.) <c><v>\greheightstar</v></c>(:) et(j) in(j) sǽ(j)cu(j)la(j) sæ(j)cu(j)<i>ló</i>(ji)<i>rum</i>.(hi) <b>A</b>(h)men.(gh..) (::)}
\paragraphspace
\gresetinitiallines{1}\greannotation{2. Ant.}\greannotation{3. a2}\gabcsnippet{(c4) <c>Q</c>Uan(h+ e)do(e[ll:1]d) na(g)tus(hj) es(j'_) (,) in(j)ef(i')fa(g)bí(i')li(j)ter(h) ex(gf) Vír(g_[uh:l]h)gi(h)ne,(g.) (;) tunc(g) im(g)plé(i)tae(j) sunt(hgh___) Scri(gh)ptú(e.)rae :(e.) (:) sic(g)ut(h) plú(j')vi(i)a(jkj) in(i') vel(j)lus(h') de(h)scen(g')dí(i)sti,(h.) (;) ut(h) sal(j')vum(i) fá(j')ce(i)res(g) ge(hv_GF)nus(g) hu(gf)má(e.)num :(e.) (:) te(h_i) lau(g)dá(gf)mus(d_e) De(g)us(ghg) no(e.)ster.(e.) (::)}
\separator
\end{document}